\documentclass{article}

% ================================================================
% Idioma y codificación
% ================================================================
\usepackage[spanish]{babel}
\usepackage[T1]{fontenc}
\usepackage[utf8]{inputenc}

% ================================================================
% Gráficos y layout
% ================================================================
\usepackage{graphicx}
\graphicspath{{./}{imagenes/}{antes_despues/}{estados_microinteracciones/}}
\usepackage[colorlinks=true, allcolors=blue]{hyperref}
\usepackage[letterpaper,top=2cm,bottom=5cm,left=3cm,right=3cm,marginparwidth=1.75cm]{geometry}

% ================================================================
% Matemáticas y símbolos
% ================================================================
\usepackage{amsmath}
\usepackage{amssymb}

% ================================================================
% Tablas y color
% ================================================================
\usepackage{tabularx}
\usepackage{array}
\usepackage{colortbl}
\usepackage{xcolor}
\definecolor{headerblue}{RGB}{30,60,120}

% ================================================================
% Listas, figuras y formato
% ================================================================
\usepackage{enumitem}
\usepackage{float}
\usepackage{subcaption}
\usepackage{fancyhdr}
\usepackage{parskip}

\pagestyle{fancy}
\fancyhf{}
\setlength{\headheight}{70pt}

% Definición del encabezado
\fancyhead[C]{%
    \begin{minipage}[c]{1.5cm}
        \includegraphics[width=1.5cm]{imagenes/uta.png}
    \end{minipage}%
    \hfill
    \begin{minipage}[c]{10cm}
        \centering
        \small {\textbf{UNIVERSIDAD TÉCNICA DE AMBATO}} \\[0.1cm]
        \scriptsize \textit{FACULTAD DE INGENIERÍA EN SISTEMAS ELECTRÓNICA E INDUSTRIAL} \\[0.1cm]
        \scriptsize \textbf{CARRERA SOFTWARE} \\[0.1cm]
        \scriptsize \textbf{CICLO ACADÉMICO: ENERO - JULIO 2026}
        \vspace{0.2cm}
    \end{minipage}%
    \hfill
    \begin{minipage}[c]{1.5cm}
        \centering
        \includegraphics[width=1.5cm]{imagenes/fisei.png}
    \end{minipage}%
}
\fancyfoot[C]{\thepage}

\begin{document}
\begin{center}
    {\large \textbf{INFORME DE GUÍA PRÁCTICA}} \\
\end{center}

% ================================================================
\section{Portada}
% ================================================================
\begin{flushleft}
\begin{tabular}{ l p{9cm} }
    \textbf{Tema:} & APE 4. Prototipado - Wireframes \\
    \textbf{Unidad de Organización Curricular:} & Profesional \\
    \textbf{Nivel y Paralelo:} & 5 - A \\
    \textbf{Alumnos participante(s):} & Cholota Bastidas Ariel Isai \\
                                      & García Amores Manolo José \\
                                      & Miranda Tarupi Luis Sebastián \\
                                      & Velasco Huayamave Kevin Andrés \\
    \textbf{Asignatura:} & Interacción Humano/Computador \\
    \textbf{Docente:} & Ing. Caiza Caizabuano Jose Ruben \\
\end{tabular}
\end{flushleft}

% ================================================================
\section{Informe de guía práctica}
% ================================================================

\subsection{Objetivos}
\subsubsection{General}
Crear un prototipo usable aplicando un wireframe.

\subsubsection{Específicos}
\begin{itemize}[leftmargin=1.5em, itemsep=3pt]
    \item Analizar los requisitos funcionales y relacionarlos con los usuarios y las pantallas del sistema.
    \item Diseñar el sitemap jerárquico, definiendo secciones primarias y secundarias.
    \item Construir los user flows de las tareas principales del usuario.
    \item Definir el esquema de navegación y el inventario de componentes reutilizables.
    \item Elaborar wireframes de baja fidelidad y refinarlos en un prototipo navegable evaluado con usuarios.
\end{itemize}

\subsection{Modalidad}
Presencial

\subsection{Tiempo de duración}
Presenciales: 4\\
No presenciales: 0

\subsection{Instrucciones}
\textbf{Acciones previas:}
\begin{itemize}[leftmargin=1.5em, itemsep=3pt]
    \item Ingrese al aula virtual de la asignatura en donde se halla el trabajo del tema tratado.
    \item Elabore el trabajo siguiendo las definiciones, conceptos y procesos aprendidos en clase.
    \item Desarrolle un prototipo usable aplicando software.
\end{itemize}

\textbf{Pasos:}
\begin{enumerate}[leftmargin=1.5em, itemsep=3pt]
    \item Revisar el material entregado por el docente sobre la aplicación de Wireframes en la creación de prototipos.
    \item Generar un prototipo aplicando elementos usables de interacción definidos en el software respectivo.
    \item Aplicar los criterios de un buen diseño en la definición del prototipo.
    \item Elaborar las conclusiones del trabajo.
    \item Subir el archivo a la plataforma dentro de la fecha indicada.
\end{enumerate}

\subsection{Listado de equipos, materiales y recursos}
\textbf{Listado de equipos y materiales generales empleados en la guía práctica:}
\begin{itemize}[leftmargin=1.5em, itemsep=3pt]
    \item Computadora
    \item Herramientas ofimáticas
    \item Aula virtual
    \item Internet
    \item Visual Studio Code y Figma
\end{itemize}

\textbf{TAC (Tecnologías para el Aprendizaje y Conocimiento) empleados en la guía práctica:}
\begin{itemize}[leftmargin=1.5em, itemsep=2pt]
    \item[$\square$] Plataformas educativas
    \item[$\square$] Simuladores y laboratorios virtuales
    \item[$\square$] Aplicaciones educativas
    \item[$\square$] Recursos audiovisuales
    \item[$\square$] Gamificación
    \item[$\blacksquare$] Inteligencia Artificial
    \item Otros (Especifique): \underline{\hspace{5cm}}
\end{itemize}

\subsection{Actividades por desarrollar}
Desarrollar la Guía APE en base al formato y las instrucciones brindadas:
\begin{enumerate}[leftmargin=1.5em, itemsep=3pt]
    \item \textbf{Análisis de requisitos:} identificar usuarios, necesidades y funcionalidades, y relacionarlas con las pantallas.
    \item \textbf{Diseño del sitemap:} construir el mapa del sitio con sus jerarquías.
    \item \textbf{Diseño del user flow:} representar el recorrido de las tareas principales con decisiones y errores.
    \item \textbf{Navegación y componentes:} definir el menú, los mecanismos de navegación y los componentes reutilizables.
    \item \textbf{Wireframes de baja fidelidad:} representar la distribución del contenido y los controles, y refinarlos en un prototipo navegable.
\end{enumerate}

\subsection{Resultados obtenidos}
El resultado de la práctica es la creación de un prototipo usable aplicando un wireframe, formando un desarrollador que pueda lograr una vista o idea de su proyecto de manera fácil. A continuación se detallan los artefactos elaborados para el caso didáctico \textbf{CaféU}, desde la Fase 1 (baja fidelidad) hasta la Fase 2 (prototipo navegable, evaluación con usuarios y mejora).

\subsubsection{Matriz de requisitos, usuarios y pantallas}
Se identificaron tres funcionalidades principales para la cafetería universitaria CaféU y su correspondencia con las pantallas del sistema:

\begin{table}[H]
\centering
\renewcommand{\arraystretch}{1.4}
\begin{tabularx}{\textwidth}{|X|X|p{4cm}|}
\hline
\rowcolor{headerblue}
\textcolor{white}{\textbf{Funcionalidad (Requisito)}} & \textcolor{white}{\textbf{Usuario y Necesidad}} & \textcolor{white}{\textbf{Pantallas}} \\ \hline
Consultar el menú y los productos disponibles. & El estudiante necesita ver qué hay disponible durante su descanso. & Inicio, Menú \\ \hline
Agregar un producto al pedido con personalización de cantidad. & Elegir un producto, ajustar su cantidad y agregarlo a la orden. & Detalle del producto, Mi pedido \\ \hline
Confirmar el pedido y recibir código de retiro. & Revisar y confirmar la orden para saber que fue recibida en cafetería. & Mi pedido, Confirmación \\ \hline
\end{tabularx}
\caption{Matriz de requisitos funcionales, usuarios y pantallas del caso CaféU.}
\end{table}

\subsubsection{Diagnóstico: usuario, tareas y problemas priorizados}
\textbf{Usuario identificado:} estudiante universitario que dispone de poco tiempo en su descanso y necesita consultar el menú y realizar un pedido sin confusiones.

Se definieron tres tareas principales para la validación:
\begin{enumerate}[leftmargin=1.5em, itemsep=3pt]
    \item \textbf{T1 (Localizar y seleccionar):} encontrar un producto relevante en el menú y seleccionar una acción.
    \item \textbf{T2 (Completar o modificar):} modificar la cantidad o eliminar un producto del pedido, incluyendo validación o error.
    \item \textbf{T3 (Finalizar y verificar):} confirmar el pedido y verificar que la operación fue recibida.
\end{enumerate}

La revisión de los wireframes permitió priorizar cinco problemas:
\begin{enumerate}[leftmargin=1.5em, itemsep=3pt]
    \item La acción para eliminar un producto del carrito era poco visible.
    \item El contraste de los botones secundarios y enlaces era insuficiente.
    \item No existía un estado vacío que informara cuando el carrito no tenía productos.
    \item Faltaba retroalimentación durante el procesamiento de la orden.
    \item El botón final de confirmación no destacaba lo suficiente.
\end{enumerate}

\subsubsection{Sitemap (Mapa del Sitio)}
Se estructuró jerárquicamente la arquitectura de la información partiendo desde la pantalla de \textbf{Inicio}, la cual conecta con el \textbf{Menú} y con \textbf{Mi pedido}. Desde el Menú se accede al \textbf{Detalle del producto}, y al confirmar el carrito se avanza hacia la \textbf{Confirmación}.

\subsubsection{User Flows (Flujos de Usuario)}
Se diseñaron tres flujos de trabajo principales para evitar callejones sin salida:
\begin{itemize}[leftmargin=1.5em, itemsep=3pt]
    \item \textbf{Flujo 1 (Realizar un pedido):} Inicio $\rightarrow$ Menú $\rightarrow$ Detalle del producto $\rightarrow$ Agregar $\rightarrow$ Mi pedido $\rightarrow$ Confirmar $\rightarrow$ Confirmación.
    \item \textbf{Flujo 2 (Editar o cancelar pedido):} Mi pedido $\rightarrow$ Cambiar cantidad / Quitar producto $\rightarrow$ Actualizar total $\rightarrow$ Confirmar o cancelar.
    \item \textbf{Flujo 3 (Buscar y consultar un producto):} Inicio $\rightarrow$ Menú $\rightarrow$ Buscar o filtrar por categoría $\rightarrow$ Detalle del producto.
\end{itemize}

\subsubsection{Esquema de navegación e inventario de componentes}
\begin{itemize}[leftmargin=1.5em, itemsep=3pt]
    \item \textbf{Menú Principal:} accesos superiores directos a Inicio, Menú y Mi pedido.
    \item \textbf{Mecanismos de Navegación:} enlaces de retroceso (ej: ``Menú''), botones de acción rápida (``Ver mi pedido'') y filtros por categoría (Bebidas, Comida).
    \item \textbf{Componentes Reutilizables:} encabezado unificado, tarjetas de producto, botones primarios/secundarios, controles de incremento (+/-) y tarjetas de resumen.
\end{itemize}

\subsubsection{Estados de la interfaz, feedback y microinteracciones}
La interfaz contempla los estados normal, de procesamiento, de error y de éxito, además de mecanismos de retroalimentación y microinteracciones:
\begin{itemize}[leftmargin=1.5em, itemsep=3pt]
    \item \textbf{Estado normal:} pantallas en reposo con los controles disponibles.
    \item \textbf{Estado de procesamiento:} al confirmar el pedido se muestra una indicación de carga hasta registrar la orden.
    \item \textbf{Estado de error:} si la operación no puede completarse, se muestra un mensaje con la causa y una opción para recuperarse.
    \item \textbf{Estado de éxito:} la pantalla de Confirmación informa que el pedido fue recibido y entrega el código de retiro \#CU-104.
    \item \textbf{Feedback de foco:} al interactuar con el buscador se aplica un contorno destacado en color rojo institucional, indicando dónde está el foco (Figura~\ref{fig:estado_foco}).
    \item \textbf{Microinteracción de filtro:} las categorías (Todo, Bebidas, Comida) cambian de color al seleccionarse, resaltando la opción activa y reduciendo la carga cognitiva.
    \item \textbf{Estado vacío:} cuando ``Mi pedido'' no tiene productos, se muestra una ilustración, el mensaje ``Aún no agregas productos'' y un botón para volver al menú.
\end{itemize}

En total se cubren cinco pantallas principales y dos estados adicionales (foco y vacío).

\begin{figure}[H]
    \centering
    \includegraphics[width=0.55\textwidth]{estados_microinteracciones/01_estado_focus_busqueda.png}
    \caption{Feedback de foco aplicado al buscador (Fase 2).}
    \label{fig:estado_foco}
\end{figure}

\subsubsection{Comparativa de wireframes: Antes (Fase 1) vs. Después (Fase 2)}
A continuación se presentan las 5 pantallas del sistema CaféU, mostrando la evolución desde el wireframe inicial de baja fidelidad hasta la versión refinada con mejoras en jerarquía visual, tipografía y contraste.

\begin{figure}[H]
    \centering
    \begin{subfigure}[b]{0.45\textwidth}
        \centering
        \includegraphics[width=0.75\textwidth]{antes_despues/w1.png}
        \caption{Inicio (Fase 1: Inicial)}
    \end{subfigure}
    \hfill
    \begin{subfigure}[b]{0.45\textwidth}
        \centering
        \includegraphics[width=0.75\textwidth]{antes_despues/w1_despues.jpeg}
        \caption{Inicio (Fase 2: Refinado)}
    \end{subfigure}
    \caption{Pantalla 1: Inicio --- Punto de entrada con acceso directo al menú y estado del pedido actual.}
\end{figure}

\begin{figure}[H]
    \centering
    \begin{subfigure}[b]{0.45\textwidth}
        \centering
        \includegraphics[width=0.75\textwidth]{antes_despues/w2.png}
        \caption{Menú (Fase 1: Inicial)}
    \end{subfigure}
    \hfill
    \begin{subfigure}[b]{0.45\textwidth}
        \centering
        \includegraphics[width=0.75\textwidth]{antes_despues/w2_despues.jpeg}
        \caption{Menú (Fase 2: Refinado)}
    \end{subfigure}
    \caption{Pantalla 2: Menú --- Buscador, filtros de categoría y tarjetas de comparación de productos.}
\end{figure}

\begin{figure}[H]
    \centering
    \begin{subfigure}[b]{0.45\textwidth}
        \centering
        \includegraphics[width=0.75\textwidth]{antes_despues/w3.png}
        \caption{Detalle (Fase 1: Inicial)}
    \end{subfigure}
    \hfill
    \begin{subfigure}[b]{0.45\textwidth}
        \centering
        \includegraphics[width=0.75\textwidth]{antes_despues/w3_despues.jpeg}
        \caption{Detalle (Fase 2: Refinado)}
    \end{subfigure}
    \caption{Pantalla 3: Detalle del producto --- Ajuste de cantidad y botón claro para agregar a la orden.}
\end{figure}

\begin{figure}[H]
    \centering
    \begin{subfigure}[b]{0.45\textwidth}
        \centering
        \includegraphics[width=0.75\textwidth]{antes_despues/w4.png}
        \caption{Mi Pedido (Fase 1: Inicial)}
    \end{subfigure}
    \hfill
    \begin{subfigure}[b]{0.45\textwidth}
        \centering
        \includegraphics[width=0.75\textwidth]{antes_despues/w4_despues.jpeg}
        \caption{Mi Pedido (Fase 2: Refinado)}
    \end{subfigure}
    \caption{Pantalla 4: Mi Pedido --- Revisión del desglose del subtotal y opción para corregir o confirmar.}
\end{figure}

\begin{figure}[H]
    \centering
    \begin{subfigure}[b]{0.45\textwidth}
        \centering
        \includegraphics[width=0.75\textwidth]{antes_despues/w5.png}
        \caption{Confirmación (Fase 1: Inicial)}
    \end{subfigure}
    \hfill
    \begin{subfigure}[b]{0.45\textwidth}
        \centering
        \includegraphics[width=0.75\textwidth]{antes_despues/w5_despues.jpeg}
        \caption{Confirmación (Fase 2: Refinado)}
    \end{subfigure}
    \caption{Pantalla 5: Confirmación --- Estado final con número de pedido (\#CU-104) para retiro.}
\end{figure}

\subsubsection{Prueba de usabilidad y accesibilidad}
\textbf{Protocolo:} se evaluó el prototipo con dos participantes ajenos al equipo de diseño, aplicando las tareas T1, T2 y T3, y registrando el éxito, la dificultad percibida y los comentarios de cada participante, sin orientar sus acciones.

\begin{table}[H]
\centering
\renewcommand{\arraystretch}{1.4}
\begin{tabularx}{\textwidth}{|l|X|c|c|X|}
\hline
\rowcolor{headerblue}
\textcolor{white}{\textbf{Participante}} & \textcolor{white}{\textbf{Tarea}} & \textcolor{white}{\textbf{¿Completó?}} & \textcolor{white}{\textbf{Dificultad}} & \textcolor{white}{\textbf{Observaciones}} \\ \hline
Usuario 1 & T1: Realizar un pedido & Sí & Baja & Recorrido fluido; sugirió hacer más visible el botón de confirmación. \\ \hline
Usuario 1 & T2: Editar y eliminar ítem & Sí & Media & Dudó al buscar la opción para quitar el producto. \\ \hline
Usuario 2 & T1: Realizar un pedido & Sí & Baja & Completó sin inconvenientes ni bloqueos. \\ \hline
Usuario 2 & T2: Editar y eliminar ítem & Sí & Baja & Sugirió aumentar el contraste de los botones secundarios. \\ \hline
\end{tabularx}
\caption{Registro de pruebas de usabilidad con dos participantes.}
\end{table}

\textbf{Heurísticas de Nielsen verificadas:}
\begin{enumerate}[leftmargin=1.5em, itemsep=3pt]
    \item \textbf{Visibilidad del estado del sistema:} el estado vacío, el feedback de foco y la pantalla de confirmación informan al usuario de la condición actual.
    \item \textbf{Reconocimiento antes que recuerdo:} la acción de eliminar usa un botón con ícono de papelera, reconocible de forma inmediata.
    \item \textbf{Consistencia y estándares:} los mismos componentes mantienen la misma función en todas las pantallas.
\end{enumerate}

\textbf{Condiciones de accesibilidad:}
\begin{enumerate}[leftmargin=1.5em, itemsep=3pt]
    \item \textbf{Contraste mínimo (WCAG 2.1, criterio 1.4.3):} se incrementó la relación de contraste de los textos a un ratio superior a 4.5:1.
    \item \textbf{Foco visible por teclado:} el contorno destacado del buscador permite identificar el foco y navegar por teclado.
\end{enumerate}

\subsubsection{Iteración basada en evidencia}
A partir de los hallazgos se corrigieron tres problemas y se comparó la versión inicial (Fase 1) con la refinada (Fase 2):

\begin{table}[H]
\centering
\renewcommand{\arraystretch}{1.4}
\begin{tabularx}{\textwidth}{|c|X|X|}
\hline
\rowcolor{headerblue}
\textcolor{white}{\textbf{\#}} & \textcolor{white}{\textbf{Hallazgo}} & \textcolor{white}{\textbf{Corrección aplicada}} \\ \hline
1 & Dificultad para localizar la acción de eliminar un producto. & Se sustituyó el enlace de texto por un botón destacado con ícono de papelera en rojo. \\ \hline
2 & Contraste insuficiente en botones secundarios y enlaces. & Se incrementó el contraste tipográfico por encima de 4.5:1. \\ \hline
3 & Falta de aviso cuando el carrito estaba vacío. & Se diseñó un estado vacío informativo con botón de retorno. \\ \hline
\end{tabularx}
\caption{Matriz de hallazgos y correcciones aplicadas en la Fase 2.}
\end{table}

Las correcciones se justifican por los principios HCI correspondientes (reconocimiento antes que recuerdo, accesibilidad y visibilidad del estado del sistema) y se evidencian en la comparativa ``Antes'' y ``Después''.

\subsubsection{Explicación didáctica del proceso}
La transición desde un requisito hasta un wireframe se demuestra con un requisito del caso CaféU:
\begin{itemize}[leftmargin=1.5em, itemsep=3pt]
    \item \textbf{Requisito:} el estudiante debe poder consultar el menú y confirmar un pedido.
    \item \textbf{Pantalla:} Menú, Mi pedido y Confirmación.
    \item \textbf{Sitemap:} Inicio $\rightarrow$ Menú $\rightarrow$ Detalle del producto; Mi pedido $\rightarrow$ Confirmación.
    \item \textbf{User flow:} Ingresar $\rightarrow$ Elegir producto $\rightarrow$ Agregar $\rightarrow$ Revisar $\rightarrow$ Confirmar.
    \item \textbf{Componentes:} tarjetas, buscador, botones, controles de cantidad y mensajes.
    \item \textbf{Wireframe:} boceto de las pantallas en baja fidelidad.
\end{itemize}

\subsubsection{Defensa y contribución de cada integrante}
Se realizó una demostración extremo a extremo del pedido completo: localizar un producto, agregarlo, revisar el carrito, confirmar y verificar la confirmación. Se mostró el error detectado (la opción para eliminar un producto resultaba poco visible) y su corrección (botón con ícono de papelera), verificando que el participante puede recuperarse sin abandonar el flujo.

\begin{itemize}[leftmargin=1.5em, itemsep=3pt]
    \item \textbf{Ariel Cholota:} wireframes de baja fidelidad y pruebas de usabilidad.
    \item \textbf{Manolo García:} esquema de navegación, inventario de componentes y documentación.
    \item \textbf{Luis Miranda:} user flows y estados/microinteracciones.
    \item \textbf{Kevin Velasco:} sitemap, coordinación y gestión del repositorio.
\end{itemize}

\subsubsection{Entregables y organización en GitHub}
Los entregables se publicaron en el repositorio \url{https://github.com/Sax300-cpu/HCI-2026-G2}, organizados en la carpeta \texttt{APE4-integral/}:
\begin{itemize}[leftmargin=1.5em, itemsep=3pt]
    \item \texttt{README.md}, \texttt{diagnostico.md} y \texttt{conclusiones.md}.
    \item \texttt{sitemap.png/pdf} y \texttt{user-flow.png/pdf}.
    \item \texttt{wireframes/}, \texttt{prototipo/} (enlace o código) y \texttt{evidencias/}.
    \item \texttt{pruebas/protocolo.md}, \texttt{pruebas/resultados.csv} y \texttt{pruebas/hallazgos.md}.
    \item \texttt{iteracion/antes-despues.md}.
\end{itemize}
Enlace al prototipo interactivo (Figma): \url{https://www.figma.com/make/9qmAFkWeKcifMwbmVkp5Q7/Caf%C3%A9U-Mobile-Web-App-Design}

\subsection{Habilidades blandas empleadas en la práctica}
\begin{itemize}[leftmargin=1.5em, itemsep=2pt]
    \item[$\blacksquare$] Liderazgo
    \item[$\blacksquare$] Trabajo en equipo
    \item[$\square$] Comunicación asertiva
    \item[$\square$] La empatía
    \item[$\square$] Pensamiento crítico
    \item[$\square$] Flexibilidad
    \item[$\square$] La resolución de conflictos
    \item[$\square$] Adaptabilidad
    \item[$\blacksquare$] Responsabilidad
\end{itemize}

\subsection{Conclusiones}
\begin{itemize}[leftmargin=1.5em, itemsep=4pt]
    \item El estudiante aprende a usar y aplicar wireframes en la creación de prototipos usables, mejorando su creatividad y reduciendo costes.
    \item La elaboración de wireframes permite estructurar la interfaz antes de la codificación, facilitando la detección temprana de problemas de navegación.
    \item La coherencia entre la matriz de requisitos, el sitemap y los flujos garantiza que ninguna funcionalidad crítica del sistema quede omitida.
\end{itemize}

\subsection{Recomendaciones}
\begin{itemize}[leftmargin=1.5em, itemsep=4pt]
    \item Previo a la elaboración de las tareas, el estudiante debe revisar la bibliografía. Si se manifiestan dudas sobre el tema de la PAE, se invita a los estudiantes a acudir a las tutorías académicas.
    \item Realizar pruebas continuas de navegación con usuarios externos para identificar posibles ambigüedades en las etiquetas y botones del sistema.
\end{itemize}

\subsection{Referencias}
\begin{itemize}[leftmargin=1.5em, itemsep=4pt]
    \item ISO. (2019). \textit{Ergonomics of human-system interaction --- Part 210: Human-centred design for interactive systems} (ISO Standard No. 9241-210:2019).
    \item Nielsen Norman Group. (2020). \textit{Wireframing 101}. \url{https://www.nngroup.com/articles/wireframes/}
    \item Nielsen Norman Group. (2021). \textit{Task Flows vs. User Flows}. \url{https://www.nngroup.com/articles/task-flows-user-flows/}
    \item Interaction Design Foundation. \textit{What is Information Architecture?}. \url{https://www.interaction-design.org/literature/topics/information-architecture}
\end{itemize}

\end{document}
